% ============================================================
% Notes: Fractional friction model
% Topics: fractional damping/friction, relaxation vs amplification,
%         stability properties, physical interpretation (memory)
% Source: Herrmann, Fractional Calculus: An Introduction for Physicists
% ============================================================

\section{Fractional friction model}

\subsection{Motivation: friction with memory}
Ordinary (local) linear friction is proportional to instantaneous velocity:
\begin{equation}
F_{\text{fr}}(t) = -\gamma\,\dot{x}(t).
\end{equation}
This is \textbf{local in time}.

In many media (viscoelastic, complex fluids, disordered environments),
the frictional response depends on the \textbf{history} of motion.
This is modeled by a \textbf{nonlocal} friction law:
\begin{equation}
F_{\text{fr}}(t) = -\int_0^t M(t-\tau)\,\dot{x}(\tau)\,d\tau,
\end{equation}
where $M(t)$ is a memory kernel (Volterra-type causal operator).

\subsection{Fractional friction kernel}
A common long-memory choice is a power-law kernel:
\begin{equation}
M(t)\propto t^{-\alpha},
\qquad 0<\alpha<1.
\end{equation}
This leads naturally to a \textbf{fractional derivative} of the velocity.

% ------------------------------------------------------------

\section{Fractional friction equation of motion}

\subsection{Fractional damping force}
A standard fractional friction model writes the friction force as
\begin{equation}
\boxed{
F_{\text{fr}}(t) = -\gamma_\alpha\,{}^{C}D_t^\alpha x(t),
\qquad 0<\alpha<1,
}
\end{equation}
or sometimes in terms of a fractional derivative of velocity:
\begin{equation}
F_{\text{fr}}(t) = -\gamma_\alpha\,{}^{C}D_t^\alpha \dot{x}(t).
\end{equation}
(Exact form depends on the physical convention and units.)

\subsection{Free fractional relaxation (prototype model)}
The simplest fractional friction dynamics is the fractional relaxation equation:
\begin{equation}
\boxed{
{}^{C}D_t^\alpha x(t) + \lambda\,x(t)=0,
\qquad 0<\alpha\le 1.
}
\end{equation}
For $\alpha=1$ this reduces to ordinary exponential relaxation:
\[
\dot{x}+\lambda x=0.
\]

\subsection{Solution: Mittag--Leffler relaxation}
The solution is expressed using the Mittag--Leffler function:
\begin{equation}
\boxed{
x(t)=x(0)\,E_\alpha(-\lambda t^\alpha),
}
\end{equation}
where
\begin{equation}
E_\alpha(z)=\sum_{k=0}^\infty \frac{z^k}{\Gamma(\alpha k+1)}.
\end{equation}

\paragraph{Key property.}
For $0<\alpha<1$, relaxation is typically \textbf{slower than exponential}
(power-law tail behavior at late times), reflecting memory.

% ------------------------------------------------------------

\section{Relaxation vs amplification}

\subsection{Relaxation (stable decay)}
Fractional relaxation corresponds to a negative ``effective rate'' in the argument:
\begin{equation}
{}^{C}D_t^\alpha x(t) + \lambda x(t)=0,
\qquad \lambda>0.
\end{equation}
Solution:
\[
x(t)=x(0)\,E_\alpha(-\lambda t^\alpha).
\]
For $\lambda>0$ this decays to $0$ (stable).

\subsection{Amplification (growth / instability)}
If the sign is reversed,
\begin{equation}
\boxed{
{}^{C}D_t^\alpha x(t) - \lambda x(t)=0,
\qquad \lambda>0,
}
\end{equation}
then
\begin{equation}
x(t)=x(0)\,E_\alpha(+\lambda t^\alpha),
\end{equation}
which grows in time (unstable amplification).

\subsection{Physical interpretation}
\begin{itemize}
\item \textbf{Relaxation:} friction/memory dissipates perturbations.
\item \textbf{Amplification:} energy is effectively injected or feedback dominates.
\end{itemize}

% ------------------------------------------------------------

\section{Main properties of fractional friction models}

\subsection{Nonlocality (memory)}
Fractional derivatives depend on the entire past:
\begin{equation}
{}^{C}D_t^\alpha x(t)
=
\frac{1}{\Gamma(1-\alpha)}
\int_0^t (t-\tau)^{-\alpha}\dot{x}(\tau)\,d\tau.
\end{equation}
Thus fractional friction is inherently \textbf{history-dependent}.

\subsection{Exponential vs Mittag--Leffler decay}
\begin{itemize}
\item Classical damping ($\alpha=1$): $x(t)\sim e^{-\lambda t}$.
\item Fractional damping ($0<\alpha<1$): $x(t)\sim E_\alpha(-\lambda t^\alpha)$
(decay is typically slower, often with algebraic tail).
\end{itemize}

\subsection{Frequency response (qualitative)}
In Fourier/Laplace space, fractional differentiation behaves like a power:
\begin{equation}
{}^{C}D_t^\alpha \ \sim\ s^\alpha.
\end{equation}
Hence fractional friction produces frequency-dependent dissipation:
\begin{equation}
\text{effective damping} \propto \omega^\alpha.
\end{equation}
This matches many complex media where damping is not constant in frequency.

\subsection{Interpolation between elastic and viscous behavior}
Fractional friction often interpolates between:
\begin{itemize}
\item purely elastic response (instantaneous / no dissipation),
\item purely viscous response (local friction proportional to $\dot{x}$),
\end{itemize}
by introducing intermediate scaling controlled by $\alpha$.

\subsection{Initial conditions}
Using Caputo derivatives, initial conditions remain classical:
\begin{equation}
x(0)=x_0,
\qquad
\dot{x}(0)=v_0 \quad (\text{if needed}).
\end{equation}
This is one reason Caputo form is preferred in physical models.

% ------------------------------------------------------------

\section{High-yield summary}

\begin{itemize}
\item Fractional friction models damping with memory (nonlocal in time).
\item Prototype fractional relaxation equation:
\[
{}^{C}D_t^\alpha x + \lambda x=0,\quad 0<\alpha\le 1.
\]
\item Solution uses Mittag--Leffler function:
\[
x(t)=x(0)\,E_\alpha(-\lambda t^\alpha).
\]
\item Relaxation: $-\lambda$ argument $\Rightarrow$ decay (stable).
\item Amplification: $+\lambda$ argument $\Rightarrow$ growth (unstable).
\item Main features: long memory, slower-than-exponential decay, frequency-dependent damping.
\end{itemize}
